% fgib_bound.tex —— FGIB 逆正交变换界与 CEB 界的紧致性比较推导
% 记号沿用 Fischer (2020)、Kolchinsky et al. (2017/2019) 与 fgib_theory.tex；用法同 dcvib_theory.tex
% 核心结论（两个方向，各带条件）：
% (1) 变分值意义：同一编码器 q 下，CEB 可学习反向编码器把 gap 消到矩匹配最小，CEB 的界 ≤ FGIB 的界，
%     差额恰为锚点 gap —— 对齐族上即 ½(s-a)² + (z/2)(σ²-1-ln σ²)，与 fgib_theory.tex 的扫参项同一表达式；
% (2) 作用对象意义：经逆正交变换 A⁻¹（mu_head 约束正交），FGIB 的界作用在"推理表示 h + 已知旁路噪声"上；
%     CEB 的界作用在训练期噪声码 Z 上，由 DPI 严格低于推理表示 μ(X) 的残差信息——作为推理表示残差的上界，
%     FGIB 的界更紧（CEB 的界根本不是该对象的有效上界），且 FGIB 的缺口可由仅旁路参数 anchor_var→0 消去。

\newtheorem{lemma}{Lemma}
\newtheorem{proposition}{Proposition}
\newtheorem{remark}{Remark}

\textbf{Notation.}
$X$ input, $Y$ label, $K$ classes.
Both models share the encoder posterior $q(z|x)$; write
$q(z|y) := \mathbb{E}_{p(x|y)}[q(z|x)]$ for the class-conditional marginal.
CEB's backward encoder is a learnable per-class (diagonal) Gaussian
$b_\theta(z|y) = \mathcal{N}(m_\theta(y), S_\theta(y))$ (linear \texttt{prior\_net} on
one-hot $y$); FGIB's prior is the fixed anchor $r(z|y) = \mathcal{N}(a\,v_y, I)$ with
orthonormal $v_y$ and anchor variance $1$.
The variational objects (Fischer 2020, Eqs.~18--20; Alemi et al.\ 2017, Eq.~14;
Kolchinsky et al.\ 2017, Eq.~10):
\begin{align*}
R &= \mathbb{E}\big[\operatorname{KL}(q(z|x)\,\|\,m(z))\big] \ge I(X;Z)
&&\text{(VIB rate, marginal prior $m$)} \\
\mathrm{ReX} &= \mathbb{E}\big[\operatorname{KL}(q(z|x)\,\|\,b_\theta(z|y))\big]
\ge I(X;Z|Y)
&&\text{(CEB residual)} \\
\widehat I &= -\frac{1}{N}\sum_i \ln\Big[\frac{1}{N}\sum_j
e^{-\|f_i-f_j\|^2/(2\sigma^2)}\Big] \ge I(X;M)
&&\text{(NIB, non-parametric)} \\
\mathrm{ReX}_F &= \mathbb{E}\big[\operatorname{KL}(q(z|x)\,\|\,r(z|y))\big]
\ge I(X;Z|Y)
&&\text{(FGIB, fixed anchor)}
\end{align*}
Fischer's comparison of the first two: at the respective prior optima
($m = q(z)$, $b = q(z|y)$), $\mathrm{ReX} = R - I(Y;Z)$ --- the conditional prior removes
exactly the retained label information, so CEB's bound is tighter than VIB's by
$I(Y;Z)$. We ask the same question one level down: where does FGIB's fixed-anchor bound
sit relative to CEB's learnable-backward bound?

\begin{lemma}[Gap identity]\label{lem:gap}
For any conditional $r(z|y)$,
\begin{equation}
\mathbb{E}\big[\operatorname{KL}(q(z|x)\,\|\,r(z|y))\big]
= I(X;Z|Y) + \mathbb{E}_y\big[\operatorname{KL}(q(z|y)\,\|\,r(z|y))\big] .
\label{eq:gap}
\end{equation}
\end{lemma}

\begin{proof}
$\mathbb{E}[\operatorname{KL}(q(z|x)\|r(z|y))] = \mathbb{E}[\ln q(z|x) - \ln r(z|y)]$
and $I(X;Z|Y) = \mathbb{E}[\ln q(z|x) - \ln q(z|y)]$ (Fischer 2020, Eqs.~9--10), where
both expectations are over $p(x,y)q(z|x)$; subtracting,
$\mathbb{E}[\ln q(z|y) - \ln r(z|y)] = \mathbb{E}_y[\operatorname{KL}(q(z|y)\|r(z|y))]$
since the argument depends on $z, y$ only. \hfill\qed
\end{proof}

\begin{lemma}[Moment matching]\label{lem:mm}
Within the per-class Gaussian (diagonal) family, the gap of Lemma~\ref{lem:gap} is
minimized by the moment-matched Gaussian $b^*_y = \mathcal{N}(m_y, S_y)$ with
$m_y = \mathbb{E}\,q(z|y)$, $S_y = \operatorname{diag}\operatorname{Var} q(z|y)$.
\end{lemma}

\begin{proof}
$\operatorname{KL}(p\|q)$ as a function of $q$ is minimized by matching the first two
moments of $p$ (standard; for Gaussians the minimizer is unique). \hfill\qed
\end{proof}

\begin{proposition}[Variational ordering: CEB's bound is tighter in value]\label{prop:ordering}
For every fixed encoder $q(z|x)$,
$\mathrm{ReX}(\theta^*) \le \mathrm{ReX}_F$ where $\theta^*$ is CEB's optimal backward
encoder, and the excess is the \emph{anchor gap}
\begin{equation}
\Delta := \mathrm{ReX}_F - \mathrm{ReX}(\theta^*)
= \mathbb{E}_y\big[\operatorname{KL}(q(z|y)\,\|\,N(a v_y, I))\big]
- \mathbb{E}_y\big[\operatorname{KL}(q(z|y)\,\|\,N(m_y, S_y))\big] \ge 0 ,
\label{eq:anchor-gap}
\end{equation}
with equality iff the anchors equal the moment-matched marginals.
On the aligned family of fgib\_theory.tex ($q(z|y) = \mathcal{N}(s v_y, \sigma^2 I)$, so
$I(X;Z|Y) = 0$ and $b^* = q(z|y)$ exactly),
\begin{equation}
\Delta = \tfrac{1}{2}(s-a)^2 + \tfrac{z}{2}\big(\sigma^2 - 1 - \ln\sigma^2\big) ,
\label{eq:aligned-gap}
\end{equation}
i.e.\ $\mathrm{ReX}_F = \Delta$ and $\mathrm{ReX}(\theta^*) = 0$: FGIB's bound value
\emph{equals} the sweep-driving compression term of fgib\_theory.tex
(Proposition~1 there), while CEB's bound collapses to the (zero) residual.
\end{proposition}

\begin{proof}
The anchor $N(a v_y, I)$ is a member of CEB's per-class Gaussian family, so the minimum
over the family is no larger than the value at the anchor: the ordering and
Eq.~\eqref{eq:anchor-gap} follow from Lemmas~\ref{lem:gap} and~\ref{lem:mm}.
On the aligned family the marginal is itself Gaussian, so $b^* = q(z|y)$ gives gap $0$,
$I(X;Z|Y) = 0$ (the posterior is class-constant), and
$\operatorname{KL}(\mathcal{N}(s v_y, \sigma^2 I)\|\mathcal{N}(a v_y, I))
= \tfrac{1}{2}(s-a)^2 + \tfrac{z}{2}(\sigma^2-1-\ln\sigma^2)$. \hfill\qed
\end{proof}

\begin{remark}[Tightness--force trade-off]\label{rem:force}
Proposition~\ref{prop:ordering} is the price of the fixed anchor, and it is paid for a
reason: the gap is the only quantity the encoder can reduce. CEB's gap is a function of
the \emph{backward encoder alone} --- $b_\theta \to q(z|y)$ eliminates it with the encoder
held fixed, so CEB's bound can become tight while the representation stays uncompressed:
its tightness certifies nothing about compression, and the gap is an unobservable
optimization residual that tracks the forward encoder. FGIB's gap can shrink only if the
\emph{encoder} moves toward the anchors (the anchor is fixed), and on the aligned family
the bound value \emph{is} the compression certificate (Eq.~\eqref{eq:aligned-gap}): bound
tightness is monotonically equivalent to compression. This is the same fixed-anchor
mechanism that keeps the sweep force persistent (fgib\_theory.tex, Proposition~1).
\end{remark}

\begin{proposition}[The inverse orthogonal transform: the object of the bound]\label{prop:inverse}
Let the side head $A$ (FGIB's introduced linear layer) be \emph{orthogonal}, and let the
trained posterior be isotropic, $\Sigma(x) = \sigma_p^2 I$ (the KL variance term is
minimized at $\sigma^2 = \sigma_p^2$, the anchor variance, so this holds at the variance
optimum). Then $\tilde z := A^{-1} z = A^{\top} z$ gives
$q'(\tilde z|x) = \mathcal{N}(h(x), \sigma_p^2 I)$ --- the posterior is centered at the
\emph{inference representation} $h$ with isotropic side noise --- and the transformed
anchors $r'(\tilde z|y) = \mathcal{N}(a\, A^{\top} v_y, I)$ are again an orthonormal frame
at scale $a$ ($A^{\top}$ preserves the frame; the isotropic family is closed under
orthogonal maps). By KL invariance under invertible linear maps (dcvib\_theory.tex,
Lemma~1),
\begin{equation}
\mathrm{ReX}_F
= I(X;\, h + \sigma_p\,\varepsilon\,|Y)
+ \mathbb{E}_y\big[\operatorname{KL}\big(q'(\tilde z|y)\,\|\,r'(\tilde z|y)\big)\big],
\qquad \varepsilon \sim \mathcal{N}(0, I) .
\label{eq:object}
\end{equation}
The bound's object is the inference representation perturbed by \emph{known, isotropic,
side-path-only} noise, and the anchor geometry is preserved exactly under the transform.
\end{proposition}

\begin{proof}
$A^{-1} = A^{\top}$ since $A$ is orthogonal; $A^{\top}\Sigma A = \sigma_p^2 I$ for
$\Sigma = \sigma_p^2 I$; $A^{\top}A = I$ for the prior covariance; orthonormality of
$\{A^{\top}v_y\}$ is preserved by orthogonality of $A^{\top}$. Equation~\eqref{eq:object}
is Lemma~\ref{lem:gap} applied to the transformed pair. \hfill\qed
\end{proof}

\begin{proposition}[Tightness on the inference object: the sense in which FGIB's bound is tighter]\label{prop:object}
(a)~\emph{CEB}. The classifier consumes the sampled code $Z$ at training and the mean
$\mu(X)$ at evaluation (Fischer 2020, Appendix~A.1). At the variance-optimal posterior,
$\sigma(x) \equiv \sigma_p(y)$ is class-constant, so $Z = \mu(X) + \sigma_p(Y)\varepsilon$
is a noisy function of $\mu(X)$ given $Y$; by the DPI,
\begin{equation}
I(X; Z | Y) \;\le\; I(X;\, \mu(X) | Y) ,
\label{eq:dpi}
\end{equation}
strictly so whenever the noise destroys within-class detail. Hence CEB's bound
$\mathrm{ReX} \ge I(X;Z|Y)$ is \emph{not} an upper bound on the residual information of
the representation actually used at inference: its object sits strictly below the
inference residual, and the deficit $I(X;\mu(X)|Y) - I(X;Z|Y)$ is the within-class
information destroyed by the training noise --- a model output ($\sigma(x)$ is trained),
which the CE term itself pressures toward $0$ (by Jensen,
$\mathbb{E}_z[\mathrm{CE}(c(z),y)] \ge \mathrm{CE}(c(\mu),y)$), driving the deficit up.
The deficit is unobservable during training.
(b)~\emph{FGIB}. The inference path is $h$ at both training and evaluation; the side
noise of Proposition~\ref{prop:inverse} is not part of any inference path. The bound's
deficit $I(X;H|Y) - I(X;\,h+\sigma_p\varepsilon\,|Y)$ is the information destroyed by a
known isotropic perturbation, and it is removable by \emph{side-path-only} parameters:
as the anchor variance $\sigma_p^2 \to 0$ (point anchors), the variance-optimal posterior
collapses ($\sigma^2 = \sigma_p^2 \to 0$) and the bound's object converges to
$I(X;H|Y)$ itself, without changing the classifier or the inference behavior.
CEB cannot take this limit: its noise is in the training path, and removing it changes
the model (the classifier is trained on the noisy codes).
(c)~\emph{Conclusion}. In the pure variational sense (same encoder, same Gaussian
family), CEB's bound is tighter (Proposition~\ref{prop:ordering}). As an upper bound on
the residual of the representation used at inference --- the object the regularizer is
meant to control --- FGIB's bound, taken through the inverse orthogonal transform, is the
tighter (in the point-anchor limit, the only valid) one: its object is the inference
representation up to a known, removable perturbation, while CEB's object is the training
code, which lies strictly below the inference residual by an unobservable, model-driven
deficit (Eq.~\eqref{eq:dpi}).
\end{proposition}

\begin{remark}[Non-orthogonal side head (current implementation)]\label{rem:orth}
Proposition~\ref{prop:inverse} is exact because the isotropic Gaussian family is closed
under orthogonal maps; the general diagonal family is not closed under arbitrary linear
maps (dcvib\_theory.tex, Remark on family closure). With an unconstrained \texttt{mu\_head}
the transformed noise $A^{-1}\Sigma A^{-\top}$ is anisotropic and the equivalence is
approximate. Constraining the side head to be orthogonal (e.g.\ Cayley or QR
parameterization) makes the inverse transform exact and preserves both the anchor frame
and the noise isotropy --- a cheap modification that certifies the bound of
Proposition~\ref{prop:object}.
\end{remark}

\begin{remark}[Anchor-frame reading]\label{rem:frame}
With the orthonormal anchor matrix $V$ (columns $v_y$), the transform $V^{\top}$ expresses
$z$ in class coordinates: the anchors become $a\,e_y$, and the gap of
Lemma~\ref{lem:gap} decomposes into $K$ independent per-class terms. The orthogonal frame
is the geometric content of the prior; both inverse transforms ($A^{\top}$ and $V^{\top}$)
preserve it.
\end{remark}

\begin{remark}[Positioning against VIB and NIB]\label{rem:position}
VIB's rate bounds the \emph{full} mutual information $I(X;Z)$ with a marginal prior;
CEB removes the retained label information ($\mathrm{ReX} = R - I(Y;Z)$ at the optima),
and FGIB inherits the residual-level object while replacing the learnable backward
encoder with a fixed geometric prior (Proposition~\ref{prop:ordering}).
NIB's non-parametric bound $\widehat I \ge I(X;M)$ is $y$-unconditioned and tight only
when the components share a covariance and form well-separated clusters (Kolchinsky et
al.\ 2017, Eq.~10 and the discussion following Eq.~11); when the noise is small relative
to the cluster spread it saturates at $\ln N$ and its gradient vanishes (their
Section~IVA) --- loose exactly in the regime where the parametric FGIB bound is
well-behaved. The two bounds are complementary: NIB certifies the full compression
non-parametrically when the geometry is favorable; FGIB certifies the residual
parametrically with a geometry of its own choosing.
\end{remark}

\begin{remark}[Degenerate anchors]\label{rem:a0}
At $a = 0$ the anchors collapse to $N(0, I)$ for all classes and the gap loses its class
structure: $\Delta = \mathbb{E}_y[\operatorname{KL}(q(z|y)\|N(0,I))]$ --- a VIB-style
marginal gap, and the object-tightness argument of
Proposition~\ref{prop:object}(b) still holds, but the compressed endpoint becomes the
trivial collapse (fgib\_theory.tex, Remark on the endpoint).
\end{remark}
